\documentclass{article}


\usepackage{PRIMEarxiv}

\usepackage[utf8]{inputenc} % allow utf-8 input
\usepackage[T1]{fontenc}    % use 8-bit T1 fonts
\usepackage{hyperref}       % hyperlinks
\usepackage{url}            % simple URL typesetting
\usepackage{booktabs}       % professional-quality tables
\usepackage{amsfonts}       % blackboard math symbols
\usepackage{nicefrac}       % compact symbols for 1/2, etc.
\usepackage{microtype}      % microtypography
\usepackage{lipsum}
\usepackage{gensymb} % for \degree symbol
\usepackage{fancyhdr}       % header
\usepackage{graphicx}       % graphics
\graphicspath{{media/}}     % organize your images and other figures under media/ folder

%Header
\pagestyle{fancy}
\thispagestyle{empty}
\rhead{ \textit{ }} 

% Update your Headers here
\fancyhead[LO]{PiggyCast: A Stacking-Based Ensemble AI Model for Weather Prediction}
% \fancyhead[RE]{Firstauthor and Secondauthor}
% Firstauthor et al. if more than 2 - must use \documentclass[twoside]{article}



  
%% Title
\title{PiggyCast: A Stacking-Based Ensemble AI Model for Weather Prediction
%%%% Cite as
%%%% Update your official citation here when published 
\thanks{\textit{\underline{Citation}}: 
\textbf{Authors. Title. Pages.... DOI:000000/11111.}} 
}

\author{
  Josiah Kimani \\
  African Institute for Mathematical Sciences (AIMS) \\
  \texttt{josiah@aims.ac.za} \\
  %% examples of more authors
   \And
  Oliver Angel{\`i}l, Chris Fotso, Jan Ravnik \\
  Ishango.ai \\
  \texttt{\{oliver@ishango.ai, chris@ishango.ai, jan@ishango.ai\}} \\
  \AND
   Koomi Amoussouvi, Benedict Ejelonu, Claire David  \\
  African Institute for Mathematical Sciences (AIMS) \\
  \texttt{\{koomi@aims.ac.za, benedict@aims.ac.za, claired@aims.ac.za\}} \\
  \And
  Ana Lučić \\
  University of Amsterdam \\
  \texttt{a.lucic@uva.nl} \\
  \And
  Luke Gregor \\
  ETH Zürich \\
  \texttt{luke.gregor@sdsc.ethz.ch} \\
  \And
  Laique Djeutchouang\\
  Stellenbosch University\\
  \texttt{ldjeutchouang@sun.ac.za}
}


\begin{document}
\maketitle


\begin{abstract}
% \lipsum[1]

Recently, AI Weather Prediction (AIWP) models have outperformed classical Numerical Models in a host of different weather prediction benchmarking criteria. Given the paradigm shift from numerical to machine learning models, such forecasts can be generated in seconds to minutes on a standard laptop. Forecast datasets from frontier AIWP models for the year 2020 have been made publicly available on the WeatherBench 2 website, facilitating independent analysis, evaluation, and further research. In this study, we introduce a traditional machine learning model trained on top of these forecast datasets (known as “stacking”) to predict ERA5, thereby exploiting the strengths of each base model and aiming to outperform forecasts from any base model alone. We coin our model \textbf{"PiggyCast"}, as we effectively piggyback off the work done by leading AI research teams with expertise and compute budgets for model training that are hard to match in an MSc thesis.  The improvement in PiggyCast's Root Mean Squared Error on Geopotential Height at 500 hPa pressure, relative to the base models, was notable, with an increase in performance as forecast lead time increased. Given the low compute cost of making forecasts, and that each frontier AIWP model has its strengths and weaknesses (depending on the weather variable, region of the globe, and forecast lead time), we argue that the future of the most skilful weather forecasts will likely come from machine learning stacking, by the very nature that stacking typically yields performance better than any base model alone.
\end{abstract}


% keywords can be removed
\keywords{Weather Forecasting \and Machine Learning \and Stacking \and Ensemble Learning \and AI Weather Prediction}



\section{Introduction}
Accurate weather prediction is crucial in disaster preparedness, resource management, and societal resilience. Traditionally, Numerical Weather Prediction (NWP) models, which rely on dynamical and physics-based simulations of atmospheric processes, have been the cornerstone of weather forecasting~\cite{scher2020artificial, rasp2024weatherbench2benchmarkgeneration}. Despite their proven skill, these models face computational challenges and sensitivity to initial conditions, especially when producing high-resolution and long-range forecasts~\cite{kieu2024predictability}.

Recent advances in Artificial Intelligence Weather Prediction (AIWP) leverage machine learning techniques to complement or even replace traditional NWP models. AIWP models efficiently learn complex patterns from large datasets, correct systematic biases in dynamical models, and offer substantial gains in computational efficiency~\cite{kochkov2023neuralgcm, benbouallegue2023hybrid}. However, purely data-driven approaches often lack physical interpretability and may struggle to generalise to unseen climate conditions or extreme weather events~\cite{rasp2024weatherbench2benchmarkgeneration}.

Hybrid forecasting systems have emerged to combine the strengths of both paradigms, integrating physics-based and data-driven predictions to improve forecast skill across various meteorological variables, including rainfall, temperature, and extreme weather phenomena~\cite{kochkov2023neuralgcm}. Concurrently, benchmarking initiatives like WeatherBench have standardised the evaluation of NWP, AI, and hybrid models, facilitating direct comparisons across datasets and metrics~\cite{rasp2024weatherbench2benchmarkgeneration, rasp2020weatherbench}.

Building on these developments, this study systematically analyses the error structures of leading AI-based, hybrid, and traditional NWP models, quantifying their similarities and differences. Moreover, it proposes an ensemble machine learning model trained on the forecasts from these models using a stacking approach, demonstrating that such a combined model can outperform any individual base model.

The key objectives of this study are to: (1) analyse error patterns across different forecasting paradigms to inform optimized model selection; (2) develop and evaluate a stacking ensemble of forecasts to enhance predictive accuracy; and (3) investigate the interpretability of the ensemble’s input features to improve understanding of the forecasting process.

This work contributes to advancing operational weather prediction by clarifying the complementary strengths and limitations of numerical, AI-based, and hybrid models. The stacking ensemble approach not only improves forecast skill but also provides insights that can guide future research in meteorology and machine learning.

\section{Data Description}
The primary dataset used in this study is the ERA5 global atmospheric reanalysis dataset, which serves as the ground truth. Forecasts from four prominent weather prediction models—IFS HRES, GraphCast, Pangu, and NeuralGCM—are evaluated against the ERA5 data. All data were accessed via the WeatherBench 2 (WB2) framework, which standardises and hosts atmospheric datasets in the Zarr format on Google Cloud, making it scalable and computationally efficient for large-scale geospatial analysis~\cite{rasp2020weatherbench,rasp2024weatherbench2benchmarkgeneration}.

The study focuses on a single key meteorological variable: Geopotential Height at the 500 hPa pressure level, evaluated globally over 9 lead times (from 48 to 240 hours in 24-hour intervals). The choice of this variable is guided by its longstanding importance in meteorological research and forecasting performance evaluation~\cite{rasp2020weatherbench,tian2007statistical,ams_level_of_nondivergence, alobaidy2022evaluation,bluestein1992synoptic}. Data are downsampled to a $5.625\degree$ ($64\times32$) resolution and temporally subsampled to 12-hour intervals for computational tractability. The year 2020 was selected for analysis.

% The forecast datasets for each model are aligned with the ERA5 reanalysis data to allow for direct comparison. The regridding process applies first-order conservative interpolation to ensure fair spatial consistency across models. All datasets used are in Zarr format, enabling fast parallel access and efficient storage.

\section{Weather Forecasting Models}
This section presents an overview of the four selected weather forecasting models evaluated in this study. Each model represents a distinct approach—numerical, hybrid, or AI-driven—and has demonstrated high accuracy in medium-range forecasting benchmarks.

\subsection{IFS HRES}
The Integrated Forecast System High-Resolution (IFS HRES) model is the flagship deterministic numerical weather prediction model developed by the European Centre for Medium-Range Weather Forecasts (ECMWF)~\cite{ecmwf2023ifs}. It uses physics-based equations that describe atmospheric dynamics and thermodynamics, and is run operationally every 12 hours. IFS HRES has long been a gold standard in weather forecasting due to its accuracy and reliability~\cite{rasp2024weatherbench2benchmarkgeneration}.

In the WB2 framework, IFS HRES forecasts serve as the performance baseline, offering a robust reference for assessing newer AI-based models. Forecasts are available in multiple spatial resolutions and pressure levels, with outputs interpolated and regridded for comparison~\cite{rasp2024weatherbench2benchmarkgeneration}.

\subsection{GraphCast}
GraphCast is a graph neural network-based AI model developed by Google DeepMind~\cite{lam2022graphcast}. It models the Earth's atmosphere as a graph, where nodes represent spatial grid points and edges capture spatial dependencies. GraphCast was trained on ERA5 reanalysis data and leverages spatiotemporal context to predict future atmospheric states.

The model has been operationalised and made available on WB2, where it provides competitive forecasts with significantly lower computational cost compared to traditional NWP models. GraphCast has been shown to outperform IFS HRES in multiple metrics, especially in medium-range forecasts~\cite{lam2022graphcast}.

\subsection{Pangu}
Pangu-Weather is a 3D vision transformer-based AI model developed by Huawei Cloud~\cite{bi2022pangu}. It adopts a hierarchical temporal structure to forecast atmospheric variables at different resolutions. Pangu has demonstrated state-of-the-art performance in several forecasting tasks, particularly in short-to-medium-range lead times.

In this study, we use Pangu's operational forecasts provided in WB2. The model forecasts have been preprocessed to match the study’s resolution and temporal sampling for consistent evaluation.

\subsection{NeuralGCM}
The Neural General Circulation Model (NeuralGCM) is a hybrid model that integrates physical principles with learned neural components~\cite{kochkov2023neuralgcm}. It aims to combine the strengths of physics-based NWP models with the data-driven efficiency of deep learning. NeuralGCM uses deep neural networks to learn subgrid-scale processes while preserving large-scale atmospheric dynamics via conservation laws.

As of the time of this study, NeuralGCM is among the leading hybrid models in WB2~\cite{rasp2024weatherbench2benchmarkgeneration}. Its deterministic forecasts are available at multiple resolutions and were selected for their high RMSE skill relative to IFS HRES on geopotential height at 500 hPa.


\section{Model Interdependency with Unsupervised Learning}
\label{sec:headings}


% \lipsum[4] 
% See Section \ref{sec:headings}.

\subsection{Multidimensional Scaling}


% \subsubsection{Headings: third level}
% \lipsum[6]

% \paragraph{Paragraph}
% \lipsum[7]


% \subsection{Multidimensional Scaling}
% \lipsum[5]
% \begin{equation}
% \xi _{ij}(t)=P(x_{t}=i,x_{t+1}=j|y,v,w;\theta)= {\frac {\alpha _{i}(t)a^{w_t}_{ij}\beta _{j}(t+1)b^{v_{t+1}}_{j}(y_{t+1})}{\sum _{i=1}^{N} \sum _{j=1}^{N} \alpha _{i}(t)a^{w_t}_{ij}\beta _{j}(t+1)b^{v_{t+1}}_{j}(y_{t+1})}}
% \end{equation}

% \subsubsection{Headings: third level}
% % \lipsum[6]

% \paragraph{Paragraph}
% \lipsum[7]


\subsection{Hierarchical Clustering with Dendrograms}

\section{PiggyCast with Supervised Learning}
\subsection{Architecture}
\subsection{Implementation}

\subsection{Evaluation}

\section{PiggyCast Interpretation with SHAPley Values}



% \section{Examples of citations, figures, tables, references}
% \label{sec:others}
% \lipsum[8] \cite{kour2014real,kour2014fast} and see \cite{hadash2018estimate}.

% The documentation for \verb+natbib+ may be found at
% \begin{center}
%   \url{http://mirrors.ctan.org/macros/latex/contrib/natbib/natnotes.pdf}
% \end{center}
% Of note is the command \verb+\citet+, which produces citations
% appropriate for use in inline text.  For example,
% \begin{verbatim}
%    \citet{hasselmo} investigated\dots
% \end{verbatim}
% produces
% \begin{quote}
%   Hasselmo, et al.\ (1995) investigated\dots
% \end{quote}

% \begin{center}
%   \url{https://www.ctan.org/pkg/booktabs}
% \end{center}


% \subsection{Figures}
% \lipsum[10] 
% See Figure \ref{fig:fig1}. Here is how you add footnotes. \footnote{Sample of the first footnote.}
% \lipsum[11] 

% \begin{figure}
%   \centering
%   \fbox{\rule[-.5cm]{4cm}{4cm} \rule[-.5cm]{4cm}{0cm}}
%   \caption{Sample figure caption.}
%   \label{fig:fig1}
% \end{figure}

% \subsection{Tables}
% \lipsum[12]
% See awesome Table~\ref{tab:table}.

% \begin{table}
%  \caption{Sample table title}
%   \centering
%   \begin{tabular}{lll}
%     \toprule
%     \multicolumn{2}{c}{Part}                   \\
%     \cmidrule(r){1-2}
%     Name     & Description     & Size ($\mu$m) \\
%     \midrule
%     Dendrite & Input terminal  & $\sim$100     \\
%     Axon     & Output terminal & $\sim$10      \\
%     Soma     & Cell body       & up to $10^6$  \\
%     \bottomrule
%   \end{tabular}
%   \label{tab:table}
% \end{table}

% \subsection{Lists}
% \begin{itemize}
% \item Lorem ipsum dolor sit amet
% \item consectetur adipiscing elit. 
% \item Aliquam dignissim blandit est, in dictum tortor gravida eget. In ac rutrum magna.
% \end{itemize}




\section{Conclusion}
Your conclusion here

\section*{Acknowledgments}
This was was supported in part by......

%Bibliography
\bibliographystyle{unsrt}  
\bibliography{references}  


\end{document}
